%% ═══════════════════════════════════════════════════════════════════════════
\subsection{Contributions}
\label{sec:contributions}
%% ═══════════════════════════════════════════════════════════════════════════

This research set out to determine whether the gap between declared and computed urgency can
be measured reliably enough, inside a working digital twin, to be used as a decision aid
rather than as a descriptive statistic. The evidence gathered answers that in two parts, and
they must not be run together.

As a device for deciding where to look, it works. On the campaign, the nodes the adequacy
signal flags in the six turns before the shortage becomes public are, without a single false
alarm, nodes that go on to miss a delivery commitment, and the worst-scoring of them is
flagged eight turns before its miss while its own manager reports that all is well. As a
device for attaching a calibrated probability to a named week, it does not work yet: the
Brier skill remains negative after a repair that halved it. The first of those is the
decision a coordinator actually takes and the one the architecture was designed around; the
second is a harder problem that the campaign turned from an unknown into a measured distance.

What bounds both statements is the evidence base rather than the design. Eight nodes of which
four are positive, eleven observable milestones and one scripted scenario cannot establish a
rate, however clean the pattern. The contributions below are therefore claimed as a working
instrument with a measured behaviour and a method for interrogating it, not as validated
performance figures.

Five contributions are claimed.

\paragraph{A signal that identifies the right nodes, early, from the gap alone.} Over the six
turns preceding the public onset of the crisis, hidden risk flags three nodes, all three of
which later miss a commitment at its originally committed deadline, with no false alarm on any
turn; three of the four eventual failures are caught, and the earliest and worst is flagged
eight turns ahead. The flag is raised on a node whose declared urgency is $0.001$, which is
the construction working as specified: neither the indicators alone nor the declaration alone
would have selected it, and only their divergence does. The signal weakens once the crisis is
public, which is the expected behaviour of a perception-gap measure rather than a defect, and
it scopes the instrument to early warning. Section~\ref{sec:res_targeting} gives the
measurement and states plainly that four positives cannot establish a rate.

\paragraph{An implemented adequacy architecture.} A multi-block noisy-OR real urgency model,
propagated bidirectionally across a supply chain DAG and confronted with a
consistency-gated, structurally elicited declared urgency through an asymmetric
Prospect-Theory penalty, is specified, implemented and operable. It fills the combination
that Section~\ref{sec:lit_gaps} identifies as absent from the reviewed literature. Its
decomposition is exact rather than approximate, which is what made the diagnosis below
possible from the application rather than by hand.

\paragraph{A diagnosed failure with an identified cause, repaired and re-measured.} The
forecast layer is shown to be badly miscalibrated on the campaign, with a Brier skill between
$-3.4$ and $-6.9$ against a base-rate forecaster and twenty-one confident predictions of
which none materialised under the production outcome definition. That failure is traced to a
single block: the time block computes milestone risk from a lead-time distribution while the
event engine writes only to capacity, cost and failure probability, so no path connects a
shock to the quantity the milestone probability reads. The claim rests on a count rather than
an inference: across the whole campaign the accumulated-delay field was written a non-zero
value exactly zero times, on every node and at every turn. Mid-crisis the block reads as
numerically zero at six of eight nodes.

The repair was then carried out and re-measured on the identical frozen forecasts, and it
changed the recommendation as well as the model. Making the events write a lead-time impact,
the obvious reading of the diagnosis, does not work: a quantity routed into the lead time is
scaled by the fraction of work remaining and therefore attenuates to nothing on the
milestones nearest their deadline. Lost time must be added to the completion date. With that
correction the Brier deficit roughly halves at every horizon, from $-0.553$ to $-0.216$ at
four weeks, and the discrimination rises; the skill nonetheless remains negative, so the
system moves from badly wrong to less wrong rather than to useful. Two further defects were
exposed by the first and are reported with it: the forecast treated declared progress as an
exact value, which made the milestone probability binary in $85\%$ of cases, and the campaign
ran with indicator coverage at $25.6\%$ and three of six urgency blocks inert, which is a
diagnosis about the data rather than the mathematics.

\paragraph{A measured first-exposure effect of showing a forecast to the declarant.} On the
turn the predictive analysis first became visible, all eight declarants either revised their
declared urgency downward or had it confirmed, and none revised upward, two turns before the
scenario's one critical event. Because the two arms are shown to carry numerically identical
forecasts, the effect is attributable to the display rather than to a different prediction.
The finding is reported for a population of language-model agents, and its extension to
human operators is unestablished.

\paragraph{A demonstration that the outcome definition dominates the reported calibration.}
Scoring identical forecasts against a plausible but naive outcome definition instead of the
production one moves the base rate from $5.0\%$ to $38.7\%$ and improves the apparent Brier
skill by a factor of more than five, while leaving discrimination essentially unchanged. A
third definition, which keeps the production notion of an adverse outcome but judges a
milestone against the deadline first committed to rather than the deadline it currently
carries, raises the four-week AUC from $0.485$ to $0.654$ and the skill from $-3.36$ to
$-0.55$. The three disagree by more than any modelling decision taken in this work. The same
trap was then met from the opposite direction: a post-hoc calibration layer fitted against a
target that differed from the product's improved calibration while degrading ranking, and was
deliberately not connected. A calibration figure reported without its outcome definition
carries little information, and a calibration layer fitted without checking its target can
degrade a system that works.

A sixth result, reported as an opening rather than a contribution, establishes that
individual loading docks are not observable from directly overhead because of canopy
occlusion rather than insufficient resolution, and that the reframed problem of delimiting an
equipped facade is both learnable and sufficient for capacity estimation.

%% ═══════════════════════════════════════════════════════════════════════════
\subsection{Status of the Research Hypotheses}
\label{sec:hypothesis_status}
%% ═══════════════════════════════════════════════════════════════════════════

Table~\ref{tab:hypothesis_status} gives the status of each pre-registered hypothesis across
the three runs. It is the single authoritative statement of those verdicts in this thesis.

\begin{table}[htbp]
\centering
\caption{Status of the six hypotheses across the executed runs}
\label{tab:hypothesis_status}
\small
\begin{tabularx}{\linewidth}{>{\raggedright\arraybackslash}p{2.5cm} >{\raggedright\arraybackslash}p{2.1cm} >{\raggedright\arraybackslash}p{2.1cm} X}
\toprule
\textbf{Hypothesis} & \textbf{Dry run} & \textbf{Arm A} & \textbf{Reading} \\
\midrule
H1 latency & Not confirmed & Not confirmed & 1 of 8 nodes on the dry run, 3 of 8 in arm A
against a threshold of 5. The direction is right and the level is short. \\
\midrule
H2 early detection & Not confirmed & Not confirmed & $H$ peaks at $0.070$ then $0.017$
against a threshold of $0.3$. The declared signal tracks the computed one too closely for a
gap to open. \\
\midrule
H3 criticality & Not confirmed & Not confirmed & Top-1 on 7 of 10 informative turns, $70\%$
against a threshold of $80\%$. Identical across runs, since the test does not depend on
$U_d$. Nine turns were excluded as saturated; Section~\ref{sec:res_repair} reports what a
log-survival sort key recovers from them, as an exploratory re-measurement rather than a
revised verdict. \\
\midrule
H4 predictive validity & Not confirmed & Not confirmed & Precision and recall both $0.00$ at
the threshold. Section~\ref{sec:res_mechanism} attributes this to the time block rather than
to the adequacy formulation. \\
\midrule
H5 contagion & Not confirmed & \textbf{Confirmed} & $\Delta U_d = +0.079$ on press turns
against $+0.005$ on calm turns. The mere-urgency effect reproduced at network scale. \\
\midrule
H6 hysteresis & Not applicable & Not applicable & Real urgency shows no measurable decline
over the recovery turns, so the comparison is not meaningful on this scenario. \\
\bottomrule
\end{tabularx}
\end{table}

One hypothesis of six is confirmed. Three qualifications belong with that count. H4 fails for
an identified mechanical reason rather than because the adequacy concept is unsound. That
mechanism has since been repaired (Section~\ref{sec:res_repair}), but H4 has \emph{not} been
re-decided here: the verdict is a precision and recall measured on the arm A declarations at
a pre-registered threshold, and re-running it on a model that has changed would not answer
the question the register asked. It should be re-decided on the human campaign, against the
repaired model, with the outcome definition written into the register. H6 was never testable
on this scenario, because the script does not de-escalate far enough for a decay rate to be
compared, which is a design fault in the scenario rather than a result.

%% ═══════════════════════════════════════════════════════════════════════════
\subsection{Limitations}
\label{sec:limitations}
%% ═══════════════════════════════════════════════════════════════════════════

Table~\ref{tab:limitations} lists the limitations judged material enough to state directly,
with their operational implication.

\begin{table}[htbp]
\centering
\caption{Documented limitations and their implications}
\label{tab:limitations}
\small
\begin{tabularx}{\linewidth}{>{\raggedright\arraybackslash}p{4.4cm} X}
\toprule
\textbf{Limitation} & \textbf{Implication} \\
\midrule
The time block was blind to shocks, and the repair does not make it good & Measured, not
suspected: on the campaign the calibrated events wrote nothing on the time axis, so milestone
risk could not respond to a disruption and read as numerically zero mid-crisis at six of
eight nodes. This is the cause of the H4 failure. The path has since been built and measured
(Section~\ref{sec:res_repair}); the Brier deficit halves and remains negative, and the
milestone probability is still binary in $71.7\%$ of cases. \\[3pt]
\midrule
The campaign ran on a quarter of its indicators & Coverage on the campaign database was
$25.6\%$, with 24 of 42 indicators empty on every node and only three of the six urgency
blocks contributing. A node with no lead time received a risk floor of zero through a
documented fallback, so an absence of data presented itself as a reassuring prediction. Any
calibration performed before this was found would have been fitted to that absence. \\[3pt]
\midrule
One predictive indicator class is missing rather than empty & The model reads current flow
but not the forward booking horizon, that is, how far ahead a node's capacity is already
sold. In the real shortage that signal preceded the degradation of every flow indicator by
months. The gap is general to any capacity-constrained link, not specific to
semiconductors. \\[3pt]
\midrule
Validation rests on agent declarants & The campaign was executed with eight isolated
language-model agents, not human consultants. Every statement about declaration behaviour is
a statement about that population. The human campaign, for which the protocol is frozen, has
not been run. \\[3pt]
\midrule
A single scenario, with a very small positive class & All results come from one
nineteen-turn scenario in which between two and six node-weeks are positive depending on
horizon. Discrimination claims are correspondingly weak in both directions. \\[3pt]
\midrule
No probabilistic guarantee & The framework produces probabilities, scores and confidence
intervals, not certainties, and the campaign showed those probabilities to be poorly
calibrated at the top of their range. \\[3pt]
\midrule
Score freshness follows input cadence & Without an ERP connector, a node's score is never
more current than its last submitted weekly review. Manual entry is an accepted trade-off,
not an oversight, and Section~\ref{sec:res_volume} reports one route out of it for a single
block. \\[3pt]
\midrule
Blocks are assumed independent & The noisy-OR aggregation requires it, and the campaign
produced a counter-example: two blocks were perfectly correlated on the arm A run, which the
tool flagged on its own. \\[3pt]
\midrule
AHP is limited to four criteria & Saaty's method degrades past roughly seven pairwise
compared criteria, which is why declared urgency stays on four while the larger set of KPI
blocks is weighted through Fuzzy BWM instead. \\[3pt]
\midrule
Event calibration priors are defensible, not measured & The Beta-Bernoulli prior weight and
the EMA reactivity coefficients follow standard conjugate Bayesian revision but are not
empirically fitted. \\[3pt]
\midrule
Event reversibility is conditional & Reverting a declared event is guaranteed correct only if
no other write has touched the same KPI in the meantime; otherwise the system raises an
explicit conflict rather than silently rebasing later writes. \\[3pt]
\midrule
Double counting is possible & An operator can both declare a calibrated event and manually
correct the same KPI in the same week. The interface warns but does not forbid it, since the
human operator remains the final authority. \\[3pt]
\midrule
The satellite capacity model is not connected & It reaches one block of six, covers
metropolitan France only, produces a structural capacity rather than a live occupancy, and
its trained segmentation model is not yet wired into the production capacity path. \\[3pt]
\midrule
The fill factor is not reproducible from the retained data & The production constant of
$0.80$ comes from a 122-site calibration whose detailed output was not kept. Recomputing over
the 68 sites still on disk gives a median of $0.97$. Volumetric outputs therefore carry an
unquantified bias of order twenty per cent until the corpus is recollected. Dock capacity,
which depends on the pitch constant instead, is unaffected. \\
\bottomrule
\end{tabularx}
\end{table}

%% ═══════════════════════════════════════════════════════════════════════════
\subsection{Future Work}
\label{sec:future_work}
%% ═══════════════════════════════════════════════════════════════════════════

The work divides into what must happen before the framework can be evaluated fairly, and
what becomes worth doing afterwards. Two items dominate both lists and are stated first
because every other entry is bounded by them.

\paragraph{Separate misperception from discretion; the present instrument cannot.} The second
chain (Section~\ref{sec:res_second_chain}) established that $[U_r - U_d]_+$ selects failing
nodes when declarants misjudge their own state and selects nothing when they judge it
correctly and choose not to escalate. Both are concealment operationally, and only the first
is visible to a questionnaire, because a questionnaire records what an actor perceives while
withholding is something an actor does. This is the sharpest open question the work raises,
and it is a question about \emph{what} is measured rather than how well. Two routes are
available and neither requires new theory: instrument the escalation channel itself, so that
the interval between a declarant recognising a problem and communicating it downstream becomes
observable alongside the declared severity; or read the free-text note the system already
collects at every turn, in which the declarants of the second chain stated their intention to
delay contacting their customers while marking near-maximal severity in the same submission.
The second route costs nothing in data collection and would be the natural first test, with
the caveat that a measure taken on generated text must be validated against human declarants
before it can carry weight.

\paragraph{Widen the evidence base; it remains the binding constraint.} The quantitative
claims of this thesis rest on two chains: eight nodes with four positives and eleven
observable milestones on one scripted scenario, and fifteen nodes with three positives and
twenty-two observable milestones on a second, unscripted one. That is why a precision of
$1.00$ over six turns is reported as an observation rather than a rate, why the confidence
interval on the four-week AUC spans $0.21$ to $0.87$, and why no movement in discrimination
between model versions can be attributed. The second chain was worth its cost --- it bounded a
central claim, which one chain could not have done --- and it also shows the limit of the
approach: two scenarios authored by the same hand cannot separate a property of supply chains
from a property of the author. None of this is repaired by a better model. It is repaired by
running the same instrument on more chains, longer horizons and more milestones, and three
routes are open without new science: replaying further documented disruptions on the same
harness, which costs only scenario authoring; instrumenting a live chain through the weekly
review, which the system already supports; and generating scale synthetically for the
calibration layer while keeping the real campaigns as the evaluation set. A campaign an order
of magnitude larger in positives would let the node-level result be stated as a rate with an
interval, and would let the calibration layer be trained on the product's own target, which is
the single step currently blocking a positive skill.

\paragraph{The time block is repaired, and doing it corrected the instruction.} This thesis
first recommended making the calibrated events write a lead-time impact. That was
implemented, measured and found wrong (Section~\ref{sec:res_repair}): lost time must be added
to the completion date, not multiplied into it, or the shock vanishes on exactly the
milestones a warning would serve. The generalisable form is that a parameter describing a
node's capacity to recover and a state accumulating what it has already lost share their
units and are easy to confuse, and the confusion stays invisible while the database is
incomplete. What remains on this axis is smoothing the milestone alarm, which today moves
from $0.05$ to $1.00$ in a single turn and is correct but unusable as advance warning.

\paragraph{Fill the indicators before calibrating anything else.} The campaign that produced
the negative result ran at $25.6\%$ coverage with half the urgency blocks inert. Completing
them from public figures of the shortage raised coverage to $49.1\%$, activated all six
blocks, and exposed one model bug that had been undetectable while the database was empty. No
amount of calibration compensates for a missing input, and this step is the only one on the
list that is general by construction: it tunes no constant to this scenario.

\paragraph{Rebuild the calibration layer's training label on the product's target.} The
reliability profile of Section~\ref{sec:res_armA_scored} is the shape that post-hoc
calibration addresses~\autocite{niculescu2005predicting,guo2017calibration}, and the
correction is a thin layer rather than a model change. Such a layer was trained off-scenario
on a synthetic factory, which respects the constraint that it must not be fitted to the
evaluation data, and measured: it improves the Brier skill by $30.8\%$ on its own training
distribution, and on the campaign it improves calibration while dropping the four-week AUC
from $0.721$ to $0.580$. Its label is an event at exactly $t+1$; the product displays a
missed milestone or a critical event over $(t, t+4]$. It is therefore not connected.
Connecting it requires the synthetic factory's snapshots to record milestone truth so that
the training label can be aligned with what the product shows, which is a concrete blocked
step rather than an open question.

\paragraph{Make the forecast state its own ignorance.} Given the arm B result, an advisory layer
that cannot anticipate exogenous shocks should not present a low probability in a confident
numeric format. Two changes follow directly: make the horizon salient in the display, so a
low four-week figure cannot be read as general safety, and state explicitly which risk
classes the model does not cover.

\paragraph{Re-run the comparison with the display controlled.} Arm B is not blinded. A design in
which some turns display a deliberately uninformative forecast would separate the effect of
the content from the effect of the display, and would establish whether the first-exposure
effect is a response to the numbers or to being given a number at all.

\paragraph{Run the human campaign.} The frozen protocol, the hashed data pack and the
operating scripts are ready. Only the human campaign can decide the hypotheses as they were
written, and it should now be run against the repaired model rather than against the
known-defective one, which would have wasted the participants.

Arm C, the closed loop in which declarants act on the terrain, has been executed in its
agent version on the repaired system (Section~\ref{sec:res_repair}), and its result sets the
expectation for the human run rather than replacing it. Judged as it must be, by final state
rather than by prediction accuracy, it met five milestones of eleven on their original
commitment against five of eleven for arm A: the loop changed which milestone was met, not
how many, saving one and replanning one that would have held. With eleven milestones that is
a fact of one campaign and not a result, and it is the quantity the human run should be
powered to measure.

\paragraph{Put the outcome definition in the pre-registration.} Section~\ref{sec:res_target_trap}
shows the outcome definition moving the reported skill by a factor of five. It belongs
alongside the numeric thresholds in the H1 to H6 register, where it currently is not.

\paragraph{Connect the observed capacity channel.} Three steps stand between the satellite work
and a usable input: wire the trained segmentation model into the production capacity path in
place of the trailer-span method, re-measure the no-signature rate on an unbiased draw, and
recollect the regulatory corpus cleanly to re-derive the fill factor. Beyond those, the
surface model from LiDAR is the most promising unexplored route, because the canopy step is
several metres lower than the main roof and would give both the canopy extent and the true
roof edge, correcting the 0 to 14.4~m overhang that currently displaces the apron boundary.
One governance question must also be settled rather than deferred: the detection library is
distributed under a copyleft licence, which is immaterial for internal use but becomes
determinative the moment the capability is distributed or exposed over a network.

\paragraph{Broaden the elicitation and couple the platforms.} Two directions identified but
deliberately descoped remain open: replacing single-operator AHP elicitation with a
multi-expert protocol aggregating several judgments per node, and coupling SupplyScore with
the DISCO platform it currently runs alongside rather than inside. Neither is a design
commitment, and both depend on what a repaired system shows.

%% ═══════════════════════════════════════════════════════════════════════════
\subsection{Personal Conclusion}
\label{sec:personal}
%% ═══════════════════════════════════════════════════════════════════════════

\subsubsection{Technical Skills Developed}

The project drew on four distinct areas, and the difficulty of the subject lay mostly in
making them work together rather than in any one of them.

\begin{itemize}
    \item \textbf{Multi-criteria decision analysis}: pairwise elicitation through AHP and
    Fuzzy BWM, consistency validation, and outranking aggregation through PROMETHEE~II.
    \item \textbf{Probabilistic modelling}: noisy-OR aggregation, Monte Carlo simulation,
    conjugate Bayesian revision in a sparse-data regime, and, added during the campaign
    analysis, probabilistic forecast verification through proper scoring rules and
    reliability decomposition.
    \item \textbf{Graph algorithms}: propagation over a multi-rank directed acyclic graph,
    topological ordering, shock simulation and network criticality.
    \item \textbf{Software engineering}: architecture of a multi-tenant Python system,
    hardened SQLite persistence with write-ahead logging, backup and migrations, a Dash web
    interface, and a milestone-based quality process with blocking test coverage and
    automated end-to-end tests.
\end{itemize}

The satellite work added a fifth: geospatial data engineering against public web services,
dataset construction and annotation, and the training and evaluation of a segmentation model
under a strict CPU budget.

\subsubsection{What I Learned}

The hardest problem was not mathematical. It was recognising, three separate times, that the
system was measuring itself. The dry run anchored declared urgency on real urgency by
construction; the agents in arm A read their briefings so carefully that their declaration
reproduced the computed state; and arm B closed the loop deliberately by showing the model's
own output to the declarant. Each time the symptom looked like a modelling problem and each
time the cause was circularity in the measurement. I would now design the independence of
the two signals first, before designing either of them.

The second lesson concerns what counts as a result. For a long stretch of the project I
treated a negative outcome as a failure to be fixed before reporting. The most useful
findings in this thesis are negative: the forecast layer is badly calibrated, the confident
tail is entirely false, one block is structurally blind, and showing that forecast to a
declarant is counterproductive. None of those would have been found by a project optimising
for a positive result, and each of them is actionable in a way that a marginal improvement
would not have been. The related discipline, which the annotation episode of
Section~\ref{sec:res_volume} taught bluntly, is that a confident self-report is not evidence,
whether it comes from an automated process or from oneself, and that the cheap verification
step should always come before the expensive one.

With hindsight I would change two things. I would have scored the forecast layer against
recorded outcomes far earlier, since doing so took a few hundred lines and revealed the time
block defect immediately, where months of implementation had not. And I would have written
the outcome definition into the pre-registration alongside the thresholds, having since
measured how much it moves the answer.

The skill that progressed most was not a technical one. It was the willingness to state a
result that does not flatter the work, in enough detail that someone else can act on it. The
subject of this thesis is, after all, the gap between what people declare and what is
actually true, and it would be a poor outcome to have measured that gap in a supply chain
while reproducing it in the reporting.

Professionally, the project confirmed a direction. I want to work where a model meets the
person who has to act on it, because that boundary is where the interesting failures live:
not in the mathematics, which can be checked, but in what a number does to the judgment of
the person reading it.
